\documentclass[12pt,a4paper]{report}
\usepackage[utf8]{inputenc}
\usepackage[T1]{fontenc}
\usepackage[spanish]{babel}

\begin{document}

\title{Cognia}
\author{Pavel Bespakhotnyi \\ Alla Kolchanova}
\date{\today}
\maketitle

\tableofcontents

\chapter{Introducción, problema y objetivos}

\section{Introducción}

Cognia es una plataforma tecnológica multiplataforma diseñada para
transformar la gestión, la organización pedagógica y la comunicación en
el sector de las clases particulares y las tutorías independientes. La
solución unifica en un único ecosistema coordinado a los tres actores
principales del proceso educativo: el profesor particular, el alumno y
sus progenitores o tutores legales.

El nombre del proyecto, Cognia, proviene del latín \textit{cognitio}
(conocimiento, aprendizaje, conciencia). El eslogan elegido es
\textit{«Clases estructuradas, progreso visible»}, que resume la
propuesta de valor: ofrecer al profesor una herramienta profesional de
gestión y al alumno y su familia una visibilidad clara del progreso
académico.

\section{Problema o necesidad detectada}

En el modelo tradicional de clases particulares, el profesor autónomo
gestiona de forma simultánea entre 10 y 25 alumnos. La ausencia de un
software especializado obliga al docente a recurrir a herramientas
genéricas e inconexas: hojas de cálculo aisladas (Excel o Google Sheets)
para los calendarios, chats de WhatsApp para avisos y recepción de
deberes, e impresos o notas en libreta para el control financiero.

Esta fragmentación operativa genera tres problemas concretos:

\begin{itemize}
    \item \textbf{Para el profesor:} dedica entre 5 y 10 horas semanales
    a tareas administrativas (cuadrar horarios, controlar cobros, enviar
    tareas, redactar informes), tiempo que resta a la preparación de sus
    contenidos didácticos.
    \item \textbf{Para el alumno:} pierde el seguimiento de sus deberes
    entre cientos de mensajes, no tiene visibilidad de sus entregas
    pendientes ni de su evolución, y pierde la motivación.
    \item \textbf{Para el padre o tutor legal:} sufre una opacidad total
    sobre el rendimiento real de su hijo y se entera de los impagos o de
    las deudas acumuladas con meses de retraso.
\end{itemize}

\subsection*{Situaciones reales de conflicto}

\begin{itemize}
    \item \textbf{Caso A (conflicto financiero):} el padre de un alumno
    recibe una notificación verbal de impago acumulado de dos meses
    (240~€) porque el profesor olvidó anotar tres sesiones dictadas
    durante días festivos y no emitió un recordatorio formal a tiempo.
    \item \textbf{Caso B (pérdida de material):} una estudiante de 3.º de
    ESO no entrega su trabajo de preparación de examen porque el enlace
    con la ficha PDF se perdió en un hilo de WhatsApp con más de 200
    mensajes cruzados.
    \item \textbf{Caso C (saturación administrativa):} un profesor
    invierte la noche del domingo (más de 3 horas) redactando
    manualmente 15 correos individuales a los padres para informar sobre
    las notas del examen semanal.
\end{itemize}

\subsection*{Limitaciones de las soluciones actuales}

Las herramientas existentes fallan al aplicarse al contexto del tutor
privado independiente:

\begin{itemize}
    \item \textbf{Aplicaciones de mensajería (WhatsApp, Telegram):}
    mezclan la vida personal con la profesional. La información académica
    se diluye en hilos informales, los archivos adjuntos caducan y no
    existe estructura de seguimiento.
    \item \textbf{Hojas de cálculo (Excel, Sheets):} requieren
    actualización manual constante, sufren errores humanos y no ofrecen
    una interfaz interactiva ni notificaciones automáticas.
    \item \textbf{Plataformas escolares institucionales (Google
    Classroom, Moodle, Alexia):} están diseñadas para colegios e
    institutos complejos. Son hipertrofiadas, requieren configuraciones
    de dominio institucional y no contemplan la relación comercial ni los
    cobros del tutor privado.
\end{itemize}

\section{Usuario o cliente objetivo}

Cognia define tres perfiles de usuario perfectamente segregados en roles,
competencias técnicas y necesidades operativas:

\begin{itemize}
    \item \textbf{Profesor particular (ROLE\_TUTOR):} autónomo que
    gestiona entre 5 y 20 alumnos. Edad: 25--50 años. Nivel técnico:
    básico--medio. Utiliza un panel web integral desde su ordenador
    portátil al finalizar su jornada académica.
    \item \textbf{Alumno (ROLE\_STUDENT):} estudiante de entre 10 y 18
    años. Nivel técnico: alto (nativo digital). Utiliza una aplicación
    móvil Android desde casa, en el centro educativo o en movimiento.
    \item \textbf{Padre o tutor legal (ROLE\_PARENT):} adulto de entre 35
    y 55 años. Nivel técnico: básico. Consulta el historial académico de
    su hijo desde su dispositivo móvil o desde el correo electrónico.
\end{itemize}

\subsection*{User Personas representativas}

\begin{itemize}
    \item \textbf{Carlos Ruiz (profesor particular):} 34 años,
    licenciado en Matemáticas, tutor privado de 18 alumnos de secundaria.
    Necesita centralizar su agenda, reducir el tiempo dedicado a cobros y
    enviar informes profesionales a los padres sin perder horas en
    correos manuales.
    \item \textbf{Lucía Gómez (alumna de 3.º de ESO):} 15 años,
    estudiante que requiere refuerzo en Física y Química. Necesita ver
    claramente sus deberes pendientes, descargar las fichas de ejercicios
    en su móvil y sentirse motivada al ver su progreso.
    \item \textbf{Roberto Gómez (padre de Lucía):} 48 años, ingeniero
    industrial con jornada laboral intensa. Necesita tener visibilidad
    del avance académico de su hija y recibir avisos claros sobre los
    pagos mensuales sin tener que preguntar constantemente.
\end{itemize}

\section{Objetivos generales del proyecto}

Los objetivos de Cognia se han formulado bajo la metodología SMART
(específicos, medibles, alcanzables, relevantes y temporales),
garantizando su verificación cuantitativa al finalizar el proyecto:

\begin{enumerate}
    \item \textbf{Optimizar el tiempo administrativo del profesor.}
    Métrica: tiempo total por sesión (registro de clase, tarea e
    informe). Criterio de éxito: inferior a 2 minutos, lo que supone una
    reducción del 80\,\% respecto al método manual.

    \item \textbf{Facilitar el acceso y la descarga de tareas en la
    aplicación Android.} Métrica: pruebas de rendimiento con Android
    Profiler. Criterio de éxito: visualización de tareas y descarga del
    PDF en menos de 1,5 segundos.

    \item \textbf{Automatizar las comunicaciones con las familias.}
    Métrica: notificaciones y correos enviados. Criterio de éxito: el
    100\,\% de los informes se envían en menos de 30 segundos tras el
    registro de la clase.

    \item \textbf{Garantizar la disponibilidad y el rendimiento de la
    API REST.} Métrica: pruebas de carga con Postman. Criterio de éxito:
    disponibilidad superior al 99\,\% y tiempo de respuesta inferior a
    300~ms en el 95\,\% de las peticiones.

    \item \textbf{Asegurar el funcionamiento sincronizado del sistema
    en las plataformas web y móvil.} Métrica: pruebas de integración
    cruzada. Criterio de éxito: los tres canales de acceso (web, Android
    y correo electrónico) operan de forma sincronizada sin errores.
\end{enumerate}

\section{Justificación del proyecto}

Cognia resuelve un problema real que afecta a las tres partes
involucradas en el proceso educativo de las clases particulares. Frente a
las alternativas existentes, el proyecto aporta las siguientes ventajas:

\begin{itemize}
    \item Centraliza en una única plataforma la gestión administrativa,
    pedagógica y comunicativa del profesor autónomo.
    \item Ofrece al alumno una experiencia móvil clara y motivadora,
    con acceso directo a sus tareas, fechas límite y progreso.
    \item Proporciona a las familias transparencia total sobre el
    rendimiento académico y contable de sus hijos.
    \item Se ajusta perfectamente al plazo de 13 \textit{sprints}
    establecido para el proyecto intermodular del ciclo DAM2.
    \item No requiere infraestructura en tiempo real ni APIs externas
    complejas, lo que garantiza su viabilidad técnica y económica.
\end{itemize}

\end{document}